\documentclass{article}
\usepackage[T1]{fontenc}
\usepackage{lmodern}
\usepackage{graphicx}
\usepackage{amsmath,amssymb}
\usepackage[a4paper,margin=2cm]{geometry}
\usepackage[absolute,overlay]{textpos}
\setlength{\TPHorizModule}{1mm}
\setlength{\TPVertModule}{1mm}
\usepackage[indonesian]{babel}
\usepackage{hyperref}
\setlength{\emergencystretch}{2em}
% unit: Halaman 19

\begin{document}

% === Halaman 19 (python/native) ===

• Hingga saat ini belum ditemukan algoritma pemfaktoran bilangan 

bulat besar dalam waktu polinomial. 

• Fakta inilah yang membuat algoritma RSA dianggap masih aman 

untuk saat ini. Semakin panjang bilangan bulatnya, maka semakin 

lama waktu yang dibutuhkan untuk memfaktorkannya. 

• Namun jika komputer kuantum semakin maju dan berkembang, maka 

algoritma RSA akan terancam, sebab n dapat difaktorkan dengan 

komputasi kuantum. Algoritma berbasis komputasi kuantum yang 

dapat memfaktorkan n dalam waktu polinom adalah algoritma Shor.

19

\end{document}
